% Source: https://artofproblemsolving.com/community/c3h2685752_quadratics_help
% Stable identity: aops-topic-2685752-q01
Let $D=b^2-4ac$.

\textbf{(a)} Since $D<0$, $\sqrt D=i\sqrt{-D}=is$ with $s>0$ real. The roots are
\[\dfrac{-b\pm is}{2a}=-\dfrac{b}{2a}\pm\dfrac{s}{2a}\,i.\]
Both $-\dfrac b{2a}$ and $\dfrac s{2a}$ are real, so the roots are $u+vi$ and $u-vi$ with $u,v$ real: complex conjugates.

\textbf{(b)} The roots are $-\dfrac b{2a}\pm\dfrac1{2a}\sqrt D$. Here $p=-\dfrac b{2a}$ and $q=\dfrac1{2a}\ne0$ are rational, so the roots are $p\pm q\sqrt D$ with $\sqrt D$ irrational.

\textbf{(c)} The numerators $-b+\sqrt D$ and $-b-\sqrt D$ already have the conjugate form (same first part, opposite second part). Dividing by the same nonzero real number $2a$ divides each part by $2a$:
\[\dfrac{P\pm Q}{2a}=\dfrac P{2a}\pm\dfrac Q{2a},\]
so the first parts stay equal and the second parts stay opposite. Because $2a$ is real (rational in (b)), the new parts are still real (rational). So the conjugate structure is not changed by the denominator. $\blacksquare$
